\section{Stochastic correction from \texorpdfstring{\cite{arous2022high}}{[arous2022high]}}

\label{2305:app:match_with_GBA}
In this Appendix we compute the diffusion correction term of the deterministic dynamic using the formalism introduced in \cite{arous2022high}; we use the notation of that paper to make the comparison easier.

We are going to analyze the plain phase retrieval problem (\(k=p=1\), no spherical constraint), using a slightly different scaling as the one we introduced in the main, in accordance with the hypothesis required by \cite{arous2022high}. Our goal is to show that our formalism contains the one presented in the paper we are analyzing here. We want to learn a predictor of the form \(f(x)=(w\cdot x)^2\), using online SGD on the data
\[x\sim \mathcal N (0,I_d) \qquad y = (w^\star \cdot x)^2 + \sqrt{\Delta} \xi, \]
with both \(w\) and \(w^\star\) with \(O_d(1)\) norm, in contrast to what we used in the main where the norm of the weights were \(O_d(\sqrt{d})\).
The SGD process is written as
\[
w^{\nu+1} = w^{\nu} - \delta_d \nabla_w L(y,w),
\]
with \(\delta_d = \sfrac\gamma d\) and \(w^0\sim\mathcal N\left(0,\sfrac1 d I_d\right)\). 
We are going to consider the high-dimensional limit, namely \({d\to+\infty}\).

We omit the verification of \(\delta\)-localizability for brevity. 

The loss and its expectation are
\[\begin{split}
L(y,w) =& \frac{1}{2} \left[(w^\star \cdot x)^2 - (w \cdot x)^2 + \sqrt{\Delta} \xi\right]^2 \\
       =& \frac12 (w^\star \cdot x)^4 + \frac12 (w \cdot x)^4 + \frac\Delta2 \xi^2 + \left[(w^\star \cdot x)^2 - (w \cdot x)^2\right] \sqrt{\Delta}\xi - \left((w^\star \cdot x)(w \cdot x)\right)^2, \\
\Phi(w) \coloneqq& \mathbb E_{x,\xi} \left[L(y,w)\right] = \frac{3}{2}\norm{w^\star}^4 + \frac{3}{2}\norm{w}^4 +\frac32\Delta  - \norm{w^\star}^2\norm{w}^2 - 2(w^\star\cdot w)^2
\end{split}\]
We introduce our summary statistics \(u = (m,q) \coloneqq (w^\star \cdot w, \norm{w}^2)\), and we further assume \(\norm{w^\star}^2=\rho\):
\[ \begin{split}
    \Phi(w) =& \frac{3\norm{w^\star}^4+\Delta}{2} + \frac{3}{2}\norm{w}^4 - \norm{w^\star}^2\norm{w}^2 - 2(w^\star\cdot w)^2 \\
    \Phi(m,q) =& \frac{3\rho^2+\Delta}{2} + \frac32 q^2 - \rho q - 2m^2
\end{split}\]
We now compute the gradients needed for the dynamics. The stochastic gradient is
\[
    \nabla_w L = 2 (w \cdot x)^3 x  - 2(w \cdot x) \sqrt{\Delta}\xi x - 2(w^\star \cdot x)^2(w \cdot x) x
\]
The gradient of the population loss is
\[
    \nabla_w\Phi = 6qw-2\rho w-4mw^\star
\]
The gradients of the sufficient statistics with respect to \(w\) are
\[
    \nabla_w m = w^\star \qquad \nabla_w q = 2w,
\]
and their Hessians are
\[
    \nabla^2_w m = 0 \qquad \nabla^2_w q = 2I_d.
\]
Last ingredient we need for the dynamic equations is \(V\coloneqq \mathbb E_{x,\xi} \left[\nabla_w(L-\Phi)\otimes\nabla_w(L-\Phi)\right]\).

We are left with
\[\begin{split}
    V ={}& \mathbb E_{x,\xi} \left[\nabla_wL\otimes\nabla_wL\right]+\nabla_w\Phi\otimes\nabla_w\Phi \\
       &- \mathbb E_{x,\xi} \left[\nabla_w L\otimes\nabla_w\Phi + \nabla_w\Phi\otimes\nabla_w L\right].
\end{split}\]
We compute the two terms separately
\[\begin{split}
    \nabla_w\Phi\otimes\nabla_w\Phi ={}& 4(3q-\rho)^2ww^\top \\
        &- 4(6qm-2\rho m)(w^\star w^{\top}+w w^{\star\top}) + 4\cdot4m^2 w^\star w^{\star\top}
\end{split}\]
Since \(\mathbb E_{x,\xi}\left[\nabla_wL\right] = \nabla_w\Phi\), the cross term reduces to the same quantity:
\[\begin{split}
    \mathbb E_{x,\xi}\left[\nabla_wL\right]\otimes\nabla_w\Phi
        =& (6\norm{w}^2w-4(w^\star\cdot w) w^\star - 2 \norm{w^\star}w)\otimes (6qw-2\rho w-4mw^\star) \\
        =& \nabla_w\Phi\otimes\nabla_w\Phi.
\end{split}\]
It remains to compute the second-moment term:
\[\begin{split}
    &\mathbb E_{x,\xi} \left[\nabla_wL\otimes\nabla_wL\right] = \\
        &\quad=\mathbb E_{x,\xi} \Big[4(w\cdot x )^6 xx^\top + 4 (w\cdot x)^2\Delta \xi^2 xx^\top \\
        &\qquad\qquad + 4 (w^\star\cdot x)^4 (w\cdot x)^2 xx^\top
                                -8(w\cdot x)^4(w^\star\cdot x)^2 xx^\top\Big] \\
        &\quad =4(90\norm{w}^4ww^\top+15\norm{w}^6I_d) + 4\Delta(\norm{w}^2I_d+2ww^\top) \\
        &\qquad +4\big[(3\norm{w^\star}^4\norm{w}^2+12\norm{w^\star}^2(w^\star\cdot w)^2)I_d \\
        &\quad\qquad +6\norm{w^\star}^4ww^\top
                + (12\norm{w^\star}^2\norm{w}^2+24(w^\star\cdot w)^2)w^\star (w^\star)^\top\\
        &\quad\qquad +24\norm{w^\star}^2(w^\star\cdot w)(w^\star w+ ww^\star)\big] \\
        &\qquad -8\big[(3\norm{w}^4\norm{w^\star}^2+12\norm{w}^2(w^\star\cdot w)^2)I_d \\
        &\quad\qquad +6\norm{w}^4w^\star(w^\star)^\top
                + (12\norm{w}^2\norm{w^\star}^2+24(w^\star\cdot w)^2)w w^\top \\
        &\quad\qquad +24\norm{w}^2(w^\star\cdot w)(w^\star w+ ww^\star)\big] \\
        %
        &\quad= 4(15q^3+\Delta q+3\rho^2q+12\rho m^2-6q^2\rho-24qm^2)I_d \\
        &\qquad + 4(90q^2+2\Delta+6\rho^2-24\rho q -48m^2)ww^\top \\
        &\qquad +4(12\rho q + 24m^2-12q^2)w^\star(w^\star)^\top \\
        &\qquad +4(24\rho m-48qm)(w^\star w+ ww^\star)
\end{split}\]

We finally get the expression for  the matrix \(V=\mathbb E_{x,\xi} \left[\nabla_wL\otimes\nabla_wL\right]-\nabla_w\Phi\otimes\nabla_w\Phi\)
\begin{equation}\label{2305:eq:matrixV}\begin{split}
   V ={}& 4(15q^3+\Delta q+3\rho^2q+12\rho m^2-6q^2\rho-24qm^2)I_d \\
        &+ 4(90q^2+2\Delta+6\rho^2-24\rho q -48m^2-9q^2+6\rho q-\rho^2)ww^\top \\
        &+ 4(12\rho q + 24m^2-12q^2-4m^2)w^\star(w^\star)^\top \\
        &+ 4(24\rho m-48qm+6qm-2\rho m)(w^\star w+ ww^\star)
\end{split}\end{equation}

We now have all the ingredients needed for computing the dynamics
\[\begin{split}
    f_m =& -\lim_{d\to+\infty}\nabla_w \Phi \cdot \nabla_w m = -6qm+2\rho m +4\rho m = 6m(\rho-q) \\
    f_q =& -\lim_{d\to+\infty}\nabla_w \Phi \cdot \nabla_w q = -12q^2 + 4\rho q + 8m^2 \\
    g_m =& \lim_{d\to+\infty}\frac{\delta_d}{2}\Tr{V \nabla^2_w m}= 0\\
    g_q =& \lim_{d\to+\infty}\frac{\delta_d}{2}\Tr{V \nabla^2_w q} \\
        =& 4(15q^3+\Delta q+3\rho^2q+12\rho m^2-6q^2\rho-24qm^2)
           \cdot\lim_{d\to+\infty}\frac{\gamma}{2d}\Tr{2I_d} \\
        =& 4\gamma(15q^3+\Delta q+3\rho^2q+12\rho m^2-6q^2\rho-24qm^2)
\end{split}\]
In the last one we used the fact that the only term that can survive the \(\Tr{\cdot}/d\) are the ones proportional to the identity. For the same reason \(\Sigma = 0\) in this case. The full equations read as
\begin{equation}\label{2305:eq:phaseretrievalGBA}\begin{split}
 \dif m &= 6m(\rho-q) \dif t \\
 \dif q &= 4\Big[2m^2+\rho q-3q^2 \\
         &\qquad +\gamma(15q^3+\Delta q+3\rho^2q+12\rho m^2-6q^2\rho-24qm^2)\Big] \dif t
\end{split}\end{equation}
Note that this equations are exactly the same as \eqref{2305:eq:quadraticODE_M} and \eqref{2305:eq:quadraticODE_Q}, as expected. What is not trivial is proving that the diffusion term will also be equivalent.

\subsection{Diffusion at generic point}
In principle diffusion terms are expected to appear when the sufficient statistic are zoom around a fixed point of the equations.
Imagine now we want to find the diffusion term around a generic point \((m,q) = (\bar m, \bar q)\) that is not necessarily a fixed point of the equations. We have to zoom the sufficient statistics in the usual way 
\[(\tilde m, \tilde q) = \left(\sqrt{d}(m-\bar m), \sqrt{d}(q-\bar q)\right)\].
The drift terms diverge, but diffusion term will appear; this diffusion is the first stochastic correction to the deterministic dynamic of the sufficient statistics, and can be added to the equations \eqref{2305:eq:phaseretrievalGBA}, paying attention to multiply it by \(\sfrac1{\sqrt{d}}\) to compensate the zoom factor.

We express the loss and its gradients in terms of the new sufficient statistics.
\[\Phi(\tilde m, \tilde q) = \frac{3\rho^2+\Delta}{2} + \frac{3}{2} \left(\frac{\tilde q}{\sqrt{d}}+\bar q\right)^2 - \rho \left(\frac{\tilde q}{\sqrt{d}}+\bar q\right) - 2\left(\frac{\tilde m}{\sqrt{d}}+\bar m\right)^2\]
Its gradient with respect to \(w\) is
\[
    \nabla_w\Phi = 6 \left(\frac{\tilde q}{\sqrt{d}}+\bar q\right) w-2\rho w-4\left(\frac{\tilde m}{\sqrt{d}}+\bar m\right)w^\star.
\]
The gradients of the zoomed statistics are
\[
    \nabla_w \tilde m = \sqrt{d} w^\star \qquad \nabla_w \tilde q = 2\sqrt{d}w,
\]
and their Hessians are
\[
    \nabla^2_w \tilde m = 0 \qquad \nabla^2_w \tilde q = 2\sqrt{d}I_d.
\]
The expression of matrix \(V\) is the same as the one in equation \eqref{2305:eq:matrixV} with the following substitutions
\[\begin{split}
    q = \frac{\tilde q}{\sqrt{d}} + \bar q \qquad
    m = \frac{\tilde m}{\sqrt{d}} + \bar m.
\end{split}\]
We now compute the three entries of the diffusion matrix in turn.
\[\begin{split}
 \Sigma_{\tilde m, \tilde m} =&  \lim_{d\to+\infty} \delta_d (\nabla_w{\tilde m})^\top V \nabla_w{\tilde m}  \\ 
 %
 =& \lim_{d\to+\infty}\Big[4(15q^3+\Delta q+3\rho^2q+12\rho m^2-6q^2\rho-24qm^2)\rho \\
 &\qquad\quad + 4(90q^2+2\Delta+6\rho^2-24\rho q -48m^2-9q^2+6\rho q-\rho^2)m^2 \\
 &\qquad\quad +4(12\rho q + 24m^2-12q^2-4m^2)\rho^2 \\
 &\qquad\quad +8(24\rho m-48qm+6qm-2\rho m)\rho m \Big] \\
 %
 =& 60\rho \bar q^3+4\Delta \rho \bar q+12\rho^3\bar q+48\rho^2 \bar m^2-24\rho^2 \bar q^2-96\rho \bar q\bar m^2 \\
 &324 \bar m^2 \bar q^2 + 8\Delta\bar m^2+24\rho^2\bar m^2 -96\rho \bar q \bar m^2 -192\bar m^4+24 \rho \bar q \bar m^2-4 \rho^2\bar m^2 \\
 & +48\rho^3 \bar q + 96\rho^2\bar m^2-48\rho^2\bar q^2-16\rho^2\bar m^2 \\
 & +176\rho^2 \bar m^2-336\rho \bar q \bar m^2 \\
 %
 =& 60\rho \bar q^3+4\Delta \rho \bar q+60\rho^3\bar q+324\rho^2 \bar m^2-72\rho^2 \bar q^2-504\rho \bar q\bar m^2 \\
 & + 8\Delta\bar m^2 -192\bar m^4+324\bar m^2 \bar q^2
\end{split}\]

The off-diagonal entry is
\[\begin{split}
 \Sigma_{\tilde m, \tilde q} =&  \lim_{d\to+\infty} \delta_d (\nabla_w{\tilde q})^\top V \nabla_w{\tilde m}  \\ 
 %
 =& \lim_{d\to+\infty}\Big[(15q^3+\Delta q+3\rho^2q+12\rho m^2-6q^2\rho-24qm^2)m \\
 &\qquad\quad + (90q^2+2\Delta+6\rho^2-24\rho q -48m^2-9q^2+6\rho q-\rho^2)qm \\
 &\qquad\quad +(12\rho q + 24m^2-12q^2-4m^2)m\rho \\
 &\qquad\quad +(24\rho m-48qm+6qm-2\rho m)(m^2+\rho q) \Big] \\
 %
 =& 120\bar m \bar q^3+8\Delta \bar m \bar q+24\rho^2\bar m \bar q+96 \rho \bar m^3-48\rho\bar m \bar q^2-192\bar m^3 \bar q \\
  & 720\bar m \bar q^3+16\Delta \bar m \bar q + 48\rho^2 \bar m \bar q -192\rho\bar m \bar q^2 -384\bar m^3 \bar q -72 \bar m \bar q^3 + 48 \rho \bar m \bar q^2 -8 \rho^2 \bar m \bar q \\
  & 96 \rho^2 \bar m \bar q + 192 \rho \bar m^3 -96 \rho \bar m \bar q^2 -32 \rho \bar m^3 \\
  & 192\rho \bar m^3-384 \bar m^3\bar q + 48 \bar m^3\bar q-16\rho \bar m^3 \\
  & 192\rho^2 \bar m \bar q-384 \rho \bar m \bar q^2 + 48 \rho \bar m \bar q^2 -16\rho^2 \bar m \bar q \\
 %
 =& 768\bar m \bar q^3+24\Delta \bar m \bar q+336\rho^2\bar m \bar q+432 \rho \bar m^3-624\rho\bar m \bar q^2-912\bar m^3 \bar q 
\end{split}\]

Finally, the bottom-right entry gives
\[\begin{split}
 \Sigma_{\tilde q, \tilde q} =&  \lim_{d\to+\infty} \delta_d (\nabla_w{\tilde q})^\top V \nabla_w{\tilde q}  \\ 
 %
 =& \lim_{d\to+\infty}16\Big[(15q^3+\Delta q+3\rho^2q+12\rho m^2-6q^2\rho-24qm^2)q \\
 &\qquad\quad + (90q^2+2\Delta+6\rho^2-24\rho q -48m^2-9q^2+6\rho q-\rho^2)q^2 \\
 &\qquad\quad +(12\rho q + 24m^2-12q^2-4m^2)m^2 \\
 &\qquad\quad +2(24\rho m-48qm+6qm-2\rho m)qm \Big] \\
 %
 =& 240 \bar q^4+16\Delta \bar q^2+48\rho^2\bar q^2+192\rho \bar m^2 \bar q-96\rho \bar q^3-384\bar q^2 \bar m^2 \\
  & 1440 \bar q^4 + 32\Delta \bar q^2 + 96\rho^2\bar q^2 - 384\rho \bar q^3 -768 \bar m ^2 \bar q^2 - 144 \bar q^4 +96\rho \bar q^3 - 16\rho^2 \bar q^2 \\
  & 192 \rho \bar m^2 \bar q + 384 \bar m^4 - 192 \bar m^2 \bar q^2 - 64 \bar m^4 \\
  & 768 \rho \bar m^2 \bar q - 1536 \bar m^2 \bar q^2 + 192 \bar m^2 \bar q^2 - 64 \rho \bar m^2 \bar q \\
 %
 =& 1536 \bar q^4+48\Delta \bar q^2 + 128\rho^2\bar q^2+1088\rho \bar m^2 \bar q-384\rho \bar q^3-2688\bar q^2 \bar m^2 +320 \bar m^4
\end{split}\]
The expressions we just computed match the one we obtained in Equation~\eqref{2305:eq:generic_p_variance}, if we set \(\gamma=0\). This difference arise from the \(\delta\)-localizability condition, that essentially force us to work with a vanishing learning rate. Remark 2 in \cite{arous2022high} explains how to generalize the formula for \(\Sigma\) to match exactly our expression.

% %%%%%%%%%%%%%%%%%%%%%%%%%%%%%%%%%%%%%%%%%%%%%%%%%%%%%%%%%%%%%%%%%%%%%%%%%%%%%%%
